\scanpage{117}

Soit $G$ un grpe \deleted{\unsure{pro}} alg comm sur $\mathbb Q$ provenant
d'une situation motivique \inserted{$\unsure{M|\mathcal M}$}, soit $U=\operatorname{Ker}(G^0\mathrel{\mathop{\unsure{-|\longrightarrow}}\limits^{\epsilon}}\mathbb G_m)$.
On a \unsure{donc|encore} un élément canonique dans $\mathbb G_m$.
\deleted{On a \unsure{$U$ tore|$U=$} \unsure{ssi|$G$}.}
% slanted
\margin{N.B. \unsure{Je n'utilise|La condition} que les \unsure{élém.|éléments|idées} de structure $1,2,3$}

\uline{Proposition} 1. Conditions équivalentes

(i) $U$ est un tore

(ii) $\epsilon$ est primitif dans l'ens des caractères $G^0\mathrel{-}\mathbb G_m$

(iii) $M_{\mathbb R}$ admet un foncteur fibre \inserted{``normalisé''}, i.e. $\unsure{\mathscr G|\mathcal G}_{\mathbb R}\ne\varnothing$.
% slanted
\margin{i.e. compatible \unsure{sur les|avec les|à un} motifs eff. de poids $0$ avec le foncteur fibre standard \unsure{corresp.} \unsure{canonique}.}

C'est tjrs vrai si $M$ a degrés pairs, i.e. $j(\mu_2)=0$.

\uline{Proposition} 2. Conditions équivalentes

(i) $U$ est de la forme $U^0\times\mu_2$, $U^0$ un tore

(ii) $\epsilon$ est \deleted{\unsure{un prim.}} un carré \inserted{dans $\operatorname{Hom}(G^0,\mathbb G_m)$}

(iii) $M_{\mathbb R}$ n'a pas de foncteur fibre \inserted{``normalisé''}, i.e. $\mathscr G_{\mathbb R}=\varnothing$.

Le fait que (i) équivaut à (ii) \inserted{dans prop 1 resp. 2}
résulte trivialmt de \deleted{l'\unsure{équivalence} de \unsure{ce}} l'existence de $j$ avec $\epsilon j(\lambda)=\lambda^2$,
ainsi que le fait que $(1\,\mathrm{ii})\Longleftrightarrow(\text{non }\deleted{\unsure{\mathrm{(iii)}}}\,2\,\mathrm{(ii)})$.

On a $(\mathrm{iii})\Rightarrow(\mathrm{ii})$, car plus généralement pour tt
\deleted{\unsure{fonct.}} \inserted{objet} de \deleted{\unsure{degré} \unsure{objet|degré}} \inserted{degré} impair de $M_{\mathbb R}$, le \deleted{degré}\inserted{rang} doit
être pair, comme il résulte de l'existence d'une forme alternée fondamentale. D'autre
part $(\mathrm i)\Rightarrow(\mathrm{iii})$, car si $T$ est un tore
compact sur $\mathbb R$, on a $H^{\unsure{2|1|n}}(\mathbb R,T)=0$ ! Enfin, si
$M$ est à degrés pairs i.e. $j(\mu_2)=0$, alors $j=2j'$
\deleted{\uline{Cor}} et $\epsilon j'=\mathrm{id}$, donc on a (ii).

\uline{Cor} Si $M$ est quelconque, la sous-catégorie \inserted{motivique} $M^{\mathrm{pair}}$ \deleted{\unsure{admet}}
(à qui correspond $G/\deleted{\unsure{T}}j(\mu_2)$) \deleted{\unsure{est}} admet un
$\mathbb R$-foncteur fibre.

% archive sheet number: 116
